\documentclass[10pt,a4paper]{amsart}
\usepackage[margin=19mm]{geometry}
\usepackage{amsmath,amssymb,mathtools}
\usepackage[scr=rsfs,cal=euler]{mathalfa}
\usepackage{xfrac}
\usepackage[unicode,hidelinks,bookmarks=false]{hyperref}
\newcommand{\ie}{i.e.\ }
\newcommand{\cf}{cf.\ }
\newcommand{\resp}{respectively\ }
\newcommand{\Subsection}{\subsection}
\providecommand{\nobreakdash}{\nobreak\hbox{-}}
\setlength{\emergencystretch}{2em}
\multlinegap0pt
\usepackage{amssymb}
\usepackage[shortlabels]{enumitem}
\usepackage{xcolor}
\newtheorem{assumption}{Assumption}
\newtheorem{proposition}{Proposition}
\newtheorem{lemma}{Lemma}
\newcommand{\R}{\mathbb{R}}
\newcommand{\con}{\mathbf{C}}
\title{Non-asymptotic uniform in time error bounds for new and old numerical schemes for SPDEs.}
\author{Can Huang, Michela Ottobre, Gideon Simpson}
\date{Source: arXiv:2603.18944v1 (CC BY 4.0)}
\begin{document}
\maketitle

\noindent\emph{Source and licence.} Non-asymptotic uniform in time error bounds for new and old numerical schemes for SPDEs.. Version arXiv:2603.18944v1 (CC BY 4.0), distributed by arXiv under the Creative Commons Attribution 4.0 International licence, http://creativecommons.org/licenses/by/4.0/, as recorded in the arXiv licence metadata for this exact version and shown on the abstract page https://arxiv.org/abs/2603.18944v1. Full licence notice: \texttt{licenses/arxiv-2603.18944-CC-BY-4.0.txt}.\par\smallskip

\noindent\textit{Editorial scope.} Counted unit: the article's proposition taming1 (source0.tex lines 2289--2295). The article's complete original proof of it is delivered with it (source0.tex lines 2424--2511, one \texttt{proof} environment of 88 lines, which the article places away from the statement). No mathematics is written, repaired or re-derived: the statement and the proof are the article's own lines, in the article's own notation, and no retained line is substituted or re-flowed.  Retained uncounted as the source-local context the proof needs: lines 179--187; lines 442--444; lines 501--504; lines 698--701; lines 743--759; lines 762--768; lines 789--793; lines 796--802; lines 808--813; lines 1938--1943; lines 2089--2091; lines 2149--2151; lines 2159--2162; lines 2297--2299; lines 2303--2309; lines 2312--2320.  Nothing was omitted from any retained span, so the package carries no omission record.  No retained line contains a citation, so the wrapper carries no bibliography and no bibliography entry.  No retained line contains a figure environment, an image-inclusion command or drawing code, so no image is delivered and the package declares no asset dependency.  Editorial formatting only, in addition: an AMS \texttt{amsart} preamble that restates the article's own declarations and macros used by the retained lines, namely the environments \texttt{assumption}, \texttt{proposition}, \texttt{lemma} on separate counters, together with the standard packages those lines need.  The retained environments print in this document as Assumption 1, Proposition 1, Lemma 1 in order of appearance, whereas the article prints the same items with the numbering of its own full document.  The complete native source of the article as distributed in the arXiv e-print for this exact version is bundled unchanged as \texttt{sources/196739/source0.tex}.
\begin{equation}\label{eqn:SPDEintro}
\left\{
\begin{array}{ll}
du= \left(\Delta u  + f(u)\right) dt+  dW(t), & t\geq 0, x\in \mathcal O\\
u(t,x)=0, & x \in \partial \mathcal O, t\geq 0\\
u(0,x)= u_0(x), & x \in \mathcal O\, \,.
\end{array}
\right.
\end{equation}

\begin{align}\label{def:dislap}
(-\Delta_h v_h, w_h)=(\nabla v_h, \nabla w_h), \quad \mbox{for every } v_h, w_h\in V_h.
\end{align}

    \begin{align}
    & (f(u) - f(v))(u-v) \leq K_1 |u-v|^2 \label{C21} \\
    & |f'(u)| \leq K_2 (1+ |u|^{p-1}) \label{C22}
    \end{align}

\begin{equation}\label{scheme:taming1}
(u_h^{k+1}, \psi) = (u_h^{k}, \psi) +  (\Delta u_h^{k+1}, \psi) \tau +
\frac{\tau}{1+\alpha\tau \|\nabla f(u_h^k)\|^2}(f(u_h^{k}), \psi) + (\delta W^{k+1}, \psi), \ \psi\in V_h \,.
\end{equation}

\begin{assumption}\label{assump1}
The nonlinearity  
$f$ is the Nemytskii operator associated with a real-valued $C^2$ function $\tilde{f}:\R\rightarrow \R$. The function $\tilde f$ is such that
\begin{equation}\label{assump eqn: growth of f}
    \sup_{y \in  \R} \frac{|\tilde{f}(y)|+ |\tilde{f}'(y)|+|\tilde{f}''(y)|}{1+|y|^p} < \infty
\end{equation}
for some $p\geq 2$. {\it From here on, with abuse of notation, we use $f$ for both the operator and the function.}
Moreover, $f$ satisfies
%\begin{align}
%&(\nabla f(u),\nabla u)\leq -b_0\|u\|^p\|\|\nabla u\|^2, \ \ b_0>0.
%\end{align}
\begin{align}
&(f(u),u)\leq b_1 \, ,\label{eqn:coercive drift}\\
&(\nabla f(u), \nabla u)\leq b_2\|\nabla u\|^2 \, , \label{eqn:assump on gradient of drift}
\end{align}
for every $u \in H^1$, and for some $ b_1\geq 0$ and $b_2<1$. When dealing with the scheme \eqref{scheme:taming1}, we additionally require $b_2<\min\{1, \lambda_1^2\}$, where $\lambda_1$ is the smallest non-zero eigenvalue of $-\Delta$ on $\mathcal{O}$. 
\end{assumption}

\begin{assumption}\label{assump:coercivity assumption on overall drift}
The nonlinearity and the operator $\Delta$ are  such that
\begin{equation} 
    (\Delta(u-v), (u-v) |u-v|^{m-2}) + (f(u)-f(v), (u-v) |u-v|^{m-2}) \leq -  \lambda \|u-v\|_{L^m}^m, 
\end{equation}
for some $\lambda>0$ and for every $u,v \in L^m$ for which the LHS of the above is well defined and where $m$ is either equal to  $2$ or to $p$ (with $p$ as in \eqref{assump eqn: growth of f}). 
\end{assumption}

\begin{assumption}\label{assump3}
 The initial condition $u_0$ is such that 
$
\mathbb{E}\|u_0\|_{1}^q<\infty$,  for every  $q\geq 2.$
\end{assumption}

\begin{assumption}\label{assump2}
The covariance operator $Q$ satisfies
\begin{align}
\|(-\Delta)^{\frac{1}{2}}Q^{\frac{1}{2}}\|_{HS}<C,
\end{align}
where $C$ is a constant that depends only on $Q$ and $\Delta$ and $\| \cdot \|_{HS}$ denotes Hilbert Schmidt norm.  
\end{assumption}

\begin{assumption}\label{assumption: on the covariance operator Q} The covariance operator $Q$ is such that the stochastic convolution $Z$ is well defined in $L^{\infty}$ and moreover, the following holds
$$
\sup_{t\geq 0} \mathbb E \|Z(t)\|^\ell_{L^{\infty}} < C, \quad \mbox{ for all } \ell \in \mathbb N \, 
$$
where the constant $C$ may depend on $\ell$ but does not depend on $t$. 
\end{assumption}

\begin{proposition}\label{them: uniform in time moment bound for tamed}
(Time-uniform moment bound for the tamed scheme \eqref{scheme:taming1}).  Suppose the nonlinearity $f$ satisfies \eqref{eqn:coercive drift} and \eqref{eqn:assump on gradient of drift} and furthermore that Assumption \ref{assump2} and Assumption \ref{assump3} hold. Then there exists a constant  $D_n>0$, independent of $k$ and $h$,  such that for $\tau$ small enough, the following bound holds
\begin{align}
\mathbb{E} \| u_h^{k+1}\|_1^{2n}\leq  \mathbb{E} \| u_h^{0}\|_1^{2n}+D_n \,.
\end{align}
\end{proposition}

\begin{equation}\label{eqn: f polynomial}
\|f(u)-f(v)\|\leq C \|u-v\|(1+\|u\|_\infty^{p+1}+\|v\|_\infty^{p+1}) \,.
\end{equation}

\begin{align}\label{eq:p61u2}
\sup\limits_{t\in [0,T]} \|u(t)\|_{L^{m}(\Omega; H^1)}\leq C(T)(1+\|u_0\|_{L^{m}(\Omega; H^1)}+\|u_0\|^p_{L^m(\Omega; L^{2p})}) \,.
\end{align}

\begin{align}\label{eq:fH1}
\sup\limits_{t\in[0,T]}\|f(u(t))\|_{L^{m}(\Omega;{H}^1)}\leq
C (T) (1+ \|u_0\|^p_{L^m(\Omega; {H}^1)}+ \|u_0\|^{p^2}_{L^m(\Omega; L^{2p})}) \,.
\end{align}

\begin{align}\label{aux}
\breve{u}_h^{k+1}=S_{h,\tau}^{k+1}P_hu_0+\tau\sum\limits_{j=0}^k \frac{S_{h,\tau}^{k+1-j}P_hf(u(t_j))}{1+\tau\|\nabla f(u(t_j))\|^2}+\sum\limits_{j=0}^k S_{h,\tau}^{k+1-j}P_h[\delta W^{j+1}].
\end{align}

\begin{align}\label{eq:hatu}
\|\breve{u}_h^{k+1}\|_{L^{2q}(\Omega;\dot{H}^1)}&\leq \|u_0\|_{L^{2q}(\Omega;H^1)}+\tau \sum\limits_{j=0}^k \|(-\Delta_h)^{\frac{1}{2}}S_{h,\tau}^{k+1-j}P_hf(u(t_j))\|_{L^{2q}(\Omega;L^2)}\notag\\
&+\left\|\sum\limits_{j=0}^k(-\Delta_h)^{\frac{1}{2}}S_{h,\tau}^{k+1-j}P_h[\delta W^{j+1}]\right\|_{L^{2q}(\Omega;L^2)} \notag\\
&\leq \|u_0\|_{L^{2q}(\Omega;H^1)}+C\tau\sum\limits_{j=0}^k t_{k+1-j}^{-\frac{1}{2}}\|f(u(t_j))\|_{L^{2q}(\Omega;L^2)}\notag\\
&+C\bigg(\tau\sum\limits_{j=0}^k \|(-\Delta_h)^{\frac{1}{2}}S_{h,\tau}^{k+1-j}P_hQ^{\frac{1}{2}}\|^2_{HS}\bigg)^{\frac{1}{2}}\notag\\
&\leq C(T) (\con(u_0) +\|(-\Delta)^{\frac{1}{2}}Q^{\frac{1}{2}}\|_{HS}), 
\end{align}

\begin{lemma}\label{lem:erruhat}
With the notation introduced so far and with the same assumption as in Proposition \ref{Prop: one time  step estimate taming1}, the following holds: 
\begin{align}
\|u(t_k)-\breve{u}_h^k\|_{L^q(\Omega;L^2)}
\leq C(T)\con(u_0)(h+\tau^{\frac{1}{2}}) \, ,\notag
\end{align}
for $0\leq t_k \leq T$
or\  $k\leq T/\tau$. 
\end{lemma}

\begin{proposition}\label{Prop: one time  step estimate taming1}
Let $u_h^n$ and $u(t)$ be as in \eqref{scheme:taming1} and \eqref{eqn:SPDEintro}.  Suppose  \eqref{C21} holds for some constant $K_1 \in \R$.  Let also \eqref{C22} hold, together with Assumptions \ref{assump1}, Assumption \ref{assump:coercivity assumption on overall drift} (for $m=2$), Assumption \ref{assump3}, Assumptions \ref{assump2} and Assumption \ref{assumption: on the covariance operator Q}. Then, the following error estimate holds:
\begin{equation}\label{eqn:finite time bound for our scheme}
\mathbb E \|u(T) -u_h^n \|^2 \leq C(T) \con(u_0)(h^2+\tau),
\end{equation}
  where $n=T/\tau$. 
\end{proposition}

\begin{proof}[Proof of Proposition \ref{Prop: one time  step estimate taming1}]
Let us write
$$
\|u(t_k)-u_h^k\|_{L^2(\Omega;L^2)}\leq \|u(t_k)-\breve{u}_h^k\|_{L^2(\Omega;L^2)}+\|\breve{u}_h^k-u_h^k\|_{L^2(\Omega;L^2)}, 
$$
where $\breve{u}_h^k$ is defined in \eqref{aux}. The first added on the RHS of the above has already been estimated in  Lemma \ref{lem:erruhat}, so we only need to bound the second term. 

Denote $\breve{e}^k:=u_h^k-\breve{u}_h^k$. Then $\breve{e}_k$ satisfies the equation
\begin{align}
& \breve{e}^k-\breve{e}^{k-1}= \tau \Delta_h \breve{e}^k+\frac{\tau P_hf\left(u_h^{k-1}\right)}{1+\tau\left\|\nabla f\left(u_h^{k-1}\right)\right\|^2}-\frac{\tau P_h f\left(u\left(t_{k-1}\right))\right.}{1+\tau\left\|\nabla f\left(u\left(t_{k-1}\right)\right)\right\|^2} \\
& \breve{e}_0=0 .
\end{align}
Taking inner product of both sides with $\breve{e}^k$,  recalling the definition of $\Delta_h$ in \eqref{def:dislap} and the fact that $(u,v)=(P_hu,v)$ for any $ v\in V_h$, we have 
\begin{align}\label{eq:J}
& \frac{1}{2}\|\breve{e}^k\|^2-\frac{1}{2}\|\breve{e}^{k-1}\|^2+\frac{1}{2}\|\breve{e}^k-\breve{e}^{k-1}\|^2+\tau\|\nabla \breve{e}^k\|^2\notag \\
& =\left(\frac{\tau f\left(u_h^{k-1}\right)}{1+\tau\left\|\nabla f\left(u_h^{k-1}\right)\right\|^2}-\frac{\tau f\left(u\left(t_{k-1}\right)\right.}{1+\tau\left\|\nabla f\left(u\left(t_{k-1}\right)\right)\right\|^2}, \breve{e}^k\right):=J.
\end{align}

Computing further the expression for $J$ gives
$$
\begin{aligned}
J & =\frac{\tau\left(f\left(u_h^{k-1}\right)-f\left(u\left(t_{k-1}\right)\right), \breve{e}^k\right)+\tau^2\left(\left\|\nabla f\left(u\left(t_{k-1}\right)\right)\right\|^2 f\left(u_h^{k-1}\right)-\| \nabla f\left(u_h^{k-1}\right) \|^2 f\left(u\left(t_{k-1}\right)\right), \breve{e}^k\right)}{\left(1+\tau\left\|\nabla f\left(u\left(t_{k-1}\right)\right)\right\|^2\right)\left(1+\tau\left\|\nabla f\left(u_h^{k-1}\right)\right\|^2\right)} \\
& =\frac{\tau^2\left(\left\|\nabla f\left(u\left(t_{k-1}\right)\right)\right\|^2-\left\|\nabla f\left(u_h^{k-1}\right)\right\|^2\right)\left(f\left(u\left(t_{k-1}\right)\right), \breve{e}^k\right)}{\left(1+\tau\left\|\nabla f\left(u\left(t_{k-1}\right)\right)\right\|^2\right)\left(1+\tau\left\|\nabla f\left(u_h^{k-1}\right)\right\|^2\right)}+\frac{\tau\left(f\left(u_h^{k-1}\right)-f\left(u\left(t_{k-1}\right), \breve{e}^k\right)\right.}{1+\tau\left\|\nabla f\left(u_h^{k-1}\right)\right\|^2}\\
&:=J_1+J_2.
\end{aligned}
$$
By Young's inequality
\begin{align}
\mathbb{E} J_1 & \leq C \tau^3 \mathbb{E}
\bigg[\left(\left\|\nabla f\left(u_h^{k-1}\right)\right\|^2
+\|\nabla f(u(t_{k-1}))\|^2 \right)\|f(u(t_{k-1}))\|^2\bigg]+\frac{\tau}{4} \mathbb{E}\| \breve{e}^k\|^2 \notag \\
&\leq C\tau^3\bigg(\sqrt{\mathbb{E}\|\nabla f(u_h^{k-1})\|^4}+\sqrt{\mathbb{E}\|\nabla f(u(t_{k-1}))\|^4} \bigg)\sqrt{\mathbb{E}\|f(u(t_{k-1}))\|^4}
+\frac{\tau}{4} \mathbb{E}\|\breve{e}^k\|^2.
\end{align}
Using \eqref{eq:fH1}  and  our $H^1$ stability result Proposition \ref{them: uniform in time moment bound for tamed}, we land with 
\begin{align}
\mathbb{E}J_1\leq \tau^3 C(T)\con(u_0)+\frac{\tau}{4} \mathbb{E}\|\breve{e}^k\|^2.
\end{align}
Next, let us estimate $J_2$. Using the Young's inequality and \eqref{C21}, 
\begin{align}
J_2 & =\frac{\tau\left(
f\left(u\left(t_{k-1}\right)-f\left(u_h^{k-1}\right), \breve{e}^k\right)\right.}{1+\tau\left\|\nabla f\left(u_h^{k-1}\right)\right\|^2} \\
& =\frac{\tau\left(f\left(u\left(t_{k-1}\right)-f\left(\breve{u}_h^{k-1}\right), \breve{e}^k\right)+\tau\left( f\left(\breve{u}_h^{k-1}\right)-f\left(u_h^{k-1}\right), \breve{e}^{k-1}\right)\right.}{1+\tau\left\|\nabla f\left(u_h^{k-1}\right)\right\|^2} \notag\\
&+\frac{\tau\left(f\left(\breve{u}_h^{k-1}\right)-f\left(u_h^{k-1}\right), \breve{e}^k-\breve{e}^{k-1}\right)}
{1+\tau\left\|\nabla f\left(u_h^{k-1}\right)\right\|^2}\notag\\
& \leq \tau\left[\|f(u(t_{k-1}))-f\left(\breve{u}_h^{k-1}\right)\|^2+\frac{1}{4}\| \breve{e}^k \|^2\right]+ |K_1|\tau\| \breve{e}^{k-1}\|^2+\tau^2\left\| f\left(\breve{u}_h^{k-1}\right)-f\left(u_h^{k-1}\right) \right\|^2 \notag\\
&+\frac{1}{4}\|\breve{e}^k-\breve{e}^{k-1}\|^2 .
\end{align}
From Lemma \ref{lem:erruhat}, the embedding $H^1\hookrightarrow L^\infty$  and the bounds \eqref{eqn: f polynomial}, \eqref{eq:p61u2} and \eqref{eq:hatu} we then deduce
\begin{align}\label{eqn:fuuhat}
& \tau \mathbb{E} \| f\left(u\left(t_{k-1}\right)\right)-f\left(\breve{u}_h^{k-1}\right) \|^2 \\
& \leq C \tau \mathbb{E}\left[\left(1+\left\|u\left(t_{k-1}\right)\right\|_{\infty}^{2p+2}+\|\breve{u}_h^{k-1}\|_{\infty}^{2p+2}\right)\|u\left(t_{k-1}\right)-\breve{u}_h^{k-1}\|^2\right] \notag\\
& \leq C \tau \sqrt{\mathbb{E}\left(1+\left\|u\left(t_{k-1}\right)\right\|_{\infty}^{4p+4}+\|\breve{u}_h^{k-1}\|_{\infty}^{4p+4}\right)}\sqrt{\mathbb{E}\|u\left(t_{k-1}\right)-\breve{u}^{k-1}\|^4}\notag \\
& \leq C \tau \sqrt{\mathbb{E}\left(1+\left\|u\left(t_{k-1}\right)\right\|_{1}^{4p+4}+\|\breve{u}_h^{k-1}\|_{1}^{4p+4}\right)}\sqrt{\mathbb{E}\|u\left(t_{k-1}\right)-\breve{u}^{k-1}\|^4} \notag\\
& \leq C \tau (h^2+\tau) C(T) \con(u_0) 
\end{align}
and, similarly, 
\begin{align}
&\tau^2 \mathbb{E}\|f\left(\breve{u}_h^{k-1}\right)-f\left(u_h^{k-1}\right)\|^2 \notag\\
&\leq 2\tau^2\left[\mathbb{E}\|f(\breve{u}_h^{k-1})\|^2+\mathbb{E}\|f(u_h^{k-1})\|^2\right]\notag\\
&\leq C\tau^2\left[\mathbb{E}\|\breve{u}_h^{k-1}\|_\infty^{2p}+\mathbb{E}\|u_h^{k-1}\|_\infty^{2p}\right]\notag\\
&\leq C\tau^2\left[\mathbb{E}\|\breve{u}_h^{k-1}\|_{{H}^1}^{2p}+\mathbb{E}\|u_h^{k-1}\|_{{H}^1}^{2p}\right]\notag\\
&\leq C(T) \tau^2 \con(u_0).
\end{align}
Hence, we drop high-order terms, and have
\begin{align}\label{eq:estJ}
\mathbb{E} J &\leq \tau (h^2+\tau)C(T) \con(u_0)+ \frac{\tau}{2} \mathbb{E}\|\breve{e}^k\|^2+ K_1\tau \mathbb{E}\|\breve{e}^{k-1}\|^2+\frac{1}{4} \mathbb{E}\|\breve{e}^k-\breve{e}^{k-1}\|^2 .
\end{align}
Taking expectation on both sides of \eqref{eq:J} and applying \eqref{eq:estJ}, we have
\begin{align}\label{eq:hate1}
\mathbb{E}&\|\breve{e}^k\|^2\leq A(\tau)\mathbb{E}\|\breve{e}^{k-1}\|^2+\tau (h^2+\tau)C(T) \con(u_0)\notag
\end{align}
where
$$
A(\tau):=\frac{1+2 |K_1|\tau}{1-\tau}.
$$
Using the elementary inequality $(1+x)^k\leq e^{kx}, x,k>0$, we have
\begin{align}
& A(\tau)^k\leq e^{(2|K_1|+1)T}, \notag\\
&\sum\limits_{\ell=1}^{k-1}A(\tau)^\ell=\frac{1-A(\tau)^k}{1-A(\tau)}\leq \frac{e^{(2|K_1|+1)T}}{(2|K_1|+1)\tau}.
\end{align}
Therefore,
\begin{align}\label{eq:hate2}
\mathbb{E}\|\breve{e}^k\|^2 &\leq A(\tau)^k\mathbb{E}\|\breve{e}^0\|^2+\sum\limits_{\ell=1}^{k-1} A(\tau)^\ell\left[\tau (h^2+\tau)C(T) \con(u_0)\right]\notag \\
&\leq (h^2+\tau)C(T) e^{(2|K_1|+1)T} \con(u_0) \,.\notag
\end{align}
This concludes the proof. 
\end{proof}

\end{document}
